\documentclass[11pt,a4paper]{article}
\usepackage[margin=23mm,top=24mm,bottom=25mm,headheight=15pt]{geometry}
\usepackage[T1]{fontenc}
\usepackage[utf8]{inputenc}
\usepackage{mathpazo}
\usepackage[scaled=0.92]{helvet}
\usepackage{courier}
\DeclareUnicodeCharacter{2212}{\ensuremath{-}}
\DeclareUnicodeCharacter{2265}{\ensuremath{\ge}}
\DeclareUnicodeCharacter{2264}{\ensuremath{\le}}
\DeclareUnicodeCharacter{03A9}{\ensuremath{\Omega}}
\DeclareUnicodeCharacter{221A}{\ensuremath{\surd}}
\usepackage{microtype}
\usepackage[table]{xcolor}
\definecolor{Cosmic}{HTML}{26113F}
\definecolor{Purple}{HTML}{512575}
\definecolor{Muted}{HTML}{6B6274}
\definecolor{Gold}{HTML}{BB925A}
\definecolor{Pale}{HTML}{F4F0F7}
\usepackage{graphicx}
\setkeys{Gin}{keepaspectratio}
\usepackage{longtable,booktabs,array,calc}
\usepackage{amsmath,amssymb}
\usepackage{etoolbox}
\usepackage{enumitem}
\setlist{itemsep=0.25em,topsep=0.4em,leftmargin=1.6em}
\usepackage{titlesec}
\usepackage{needspace}
\titleformat{\section}{\sffamily\bfseries\Large\color{Purple}}{}{0em}{}
\titleformat{\subsection}{\sffamily\bfseries\large\color{Purple}}{}{0em}{}
\titleformat{\subsubsection}{\sffamily\bfseries\normalsize\color{Purple}}{}{0em}{}
\titlespacing*{\section}{0pt}{1.65em}{0.65em}
\titlespacing*{\subsection}{0pt}{1.2em}{0.45em}
\usepackage{fancyhdr}
\pagestyle{fancy}
\fancyhf{}
\fancyhead[L]{\sffamily\fontsize{8}{9}\selectfont\color{Muted}AGI ALPHA FLYWHEEL}
\fancyhead[R]{\sffamily\fontsize{8}{9}\selectfont\color{Muted}THE CERTIFIED BOUNDARY EDITION}
\fancyfoot[L]{\sffamily\fontsize{8}{9}\selectfont\color{Muted}RESEARCH REVIEW · 27 OCTOBER 2026}
\fancyfoot[R]{\sffamily\small\color{Purple}\thepage}
\renewcommand{\headrulewidth}{0.25pt}
\renewcommand{\footrulewidth}{0pt}
\usepackage{caption}
\captionsetup{font=small,labelfont={bf,sf},textfont=it,skip=7pt}
\usepackage{float}
\usepackage{tikz}
\usepackage{fvextra}
\DefineVerbatimEnvironment{Highlighting}{Verbatim}{breaklines,commandchars=\\\{\},fontsize=\small}
\usepackage{xurl}
\usepackage[colorlinks=true,linkcolor=Purple,urlcolor=Purple,citecolor=Purple,breaklinks=true]{hyperref}
\hypersetup{pdftitle={AGI ALPHA
FLYWHEEL — The Certified Boundary Edition},pdfauthor={Vincent Boucher — President of MONTREAL.AI},pdfsubject={Exact covering-system certificates, conventional challenge and complete conditional economics},pdfkeywords={covering systems, exact certificate, AGI ALPHA, research review}}
\urlstyle{same}
\Urlmuskip=0mu plus 2mu
\renewcommand{\texttt}[1]{{\ttfamily\small #1}}
\setlength{\parindent}{0pt}
\setlength{\parskip}{0.45em}
\renewcommand{\baselinestretch}{1.07}
\setlength{\emergencystretch}{2em}
\setlength{\tabcolsep}{5pt}
\renewcommand{\arraystretch}{1.18}
\AtBeginEnvironment{longtable}{\small}
\providecommand{\tightlist}{\setlength{\itemsep}{0pt}\setlength{\parskip}{0pt}}
\providecommand{\passthrough}[1]{#1}
\newcommand{\pandocbounded}[1]{#1}
\newlength{\cslhangindent}
\setlength{\cslhangindent}{1.5em}
\newlength{\csllabelwidth}
\setlength{\csllabelwidth}{3em}
\newenvironment{CSLReferences}[2]{\begingroup\small\raggedright\setlength{\parskip}{0.5em}\setlength{\parindent}{0pt}}{\endgroup}
\newcommand{\CSLBlock}[1]{#1\hfill\break}
\newcommand{\CSLLeftMargin}[1]{\parbox[t]{\csllabelwidth}{#1}}
\newcommand{\CSLRightInline}[1]{\parbox[t]{\linewidth-\csllabelwidth}{#1}\break}
\newcommand{\CSLIndent}[1]{\hspace{\cslhangindent}#1}
\newcommand{\evidencebreak}{\clearpage}
\begin{document}
{\sffamily\color{Purple}\small MONTREAL.AI\hfill RESEARCH REVIEW · v1}\par
\vspace{0.5em}
{\sffamily\bfseries\LARGE AGI ALPHA FLYWHEEL}\par
{\sffamily\large\color{Purple}The Certified Boundary Edition · Executive
brief}\par
\vspace{0.5em}
{\small Vincent Boucher --- President of MONTREAL.AI}\par
\noindent\textcolor{Gold}{\rule{\linewidth}{0.7pt}}
\Needspace{7\baselineskip}

\hypertarget{a-universal-answer-available-to-the-next-researcher}{%
\subsection{A universal answer, available to the next
researcher}\label{a-universal-answer-available-to-the-next-researcher}}

An executed study has produced an exact mathematical certificate
answering Gerry Myerson's documented question about finite congruence
partitions. The strongest construction came from the independent
conventional challenge. Its proof applies to every interval length and
can be checked with ordinary Python, without a model, numerical
optimizer or trusted forecast (Myerson 2026).

\begin{samepage}

Partition \(1,\ldots,n\) into complete residue classes, each hitting the
interval at least twice. Define the \textbf{heft} \(h\) as the sum of
reciprocal moduli. The study proves

\[\boxed{\displaystyle h\ge\frac{17902560}{18258529}
=0.9805039606421744\ldots>0.98.}\]

\end{samepage}

This rules out the requested below-0.98 outcome. The proof even allows
repeated moduli. It retains complete classes, exact coverage and minimum
class size, and makes no claim about singleton classes or
two-dimensional covers. The strongest inspected prior general bound is
\(\sqrt8-2\approx0.828427\) (Emanuel 2011). The new constant is not
claimed optimal; external expert review and fuller priority assessment
remain pending.

\Needspace{7\baselineskip}

\hypertarget{what-was-constructed-and-checked}{%
\subsection{What was constructed and
checked}\label{what-was-constructed-and-checked}}

A piecewise-linear weight with 21 exact integer heights converts the
question into an auditable certificate. Two separately implemented
line-arrangement checkers agree on the maximum class weight; an analytic
tail covers all larger classes. Assigning each integer its cell integral
transfers the inequality to every finite interval length without
approximation.

The first successful approach used a degree-20 rational polynomial. It
remains independently verified, but its bound is weaker. A fresh
workstream, given the public question and ordinary tools without that
witness, found the stronger construction. This is an original explicit
result relative to inspected sources, obtained with established methods.
It does not demonstrate an exclusive algorithm, model superiority or
discovery advantage.

The package preserves unsuccessful code searches, an attractive but
invalid first polynomial, its exact counterexample, both accepted
constructions, frozen configurations, fresh tests and interface
corrections. Separate checkers and adversarial tests reduce error risk;
they do not constitute external human peer review or proof-assistant
formalization.

\Needspace{7\baselineskip}

\hypertarget{useful-science-commercial-foundation-still-conditional}{%
\subsection{Useful science; commercial foundation still
conditional}\label{useful-science-commercial-foundation-still-conditional}}

\textbf{Actual customers, revenue and receipts are zero. Complete actual
cost and profit are unknown.} No profitable enterprise or Operational
Alpha is measured.

A complete planning model evaluates reproduction, qualified review and
integration of this fixed certificate package, including one year of
support. It does not price a promise to solve future open problems.
Under explicitly assumed human-resource estimates:

\begin{longtable}[]{@{}
  >{\raggedright\arraybackslash}p{(\columnwidth - 2\tabcolsep) * \real{0.5000}}
  >{\raggedleft\arraybackslash}p{(\columnwidth - 2\tabcolsep) * \real{0.5000}}@{}}
\toprule\noalign{}
\begin{minipage}[b]{\linewidth}\raggedright
Central conditional arrangement
\end{minipage} & \begin{minipage}[b]{\linewidth}\raggedleft
Amount
\end{minipage} \\
\midrule\noalign{}
\endhead
\bottomrule\noalign{}
\endlastfoot
Four prospective sponsors at an assumed \$6,000 each & \$24,000
revenue \\
Complete project cost, including discovery once & \$21,592 \\
Full-project profit before financing and tax & \$2,408 \\
Additional unrestricted financing above \$5,000 capital & \$6,573.20 \\
\end{longtable}

The central case is profitable but underfunded under its assumptions.
Refundable deposits remain restricted. An illustrative 90-day bridge
reduces profit to \$2,072.63; no lender offer exists. The adverse
scenario loses \textbf{\$19,968.80} and cannot be described as merely
underfunded. All fees, demand and labor estimates remain unvalidated.

The supported coordination mechanism is sharing fixed review costs among
compatible beneficiaries. An equally capable conventional shared
arrangement using the same artifact ties exactly. Specialist extraction
remains \textbf{HOLD}: its additional modeled \$1,042 cost has no
demonstrated benefit over giving a competent successor the same
ordinary-library asset.

\Needspace{7\baselineskip}

\hypertarget{three-horizons-and-the-next-decisive-step}{%
\subsection{Three horizons and the next decisive
step}\label{three-horizons-and-the-next-decisive-step}}

\textbf{Near term:} a durable answer, exact checkers, clear scope and a
preserved correction history. \textbf{Expansion:} fresh problems,
qualified review, paid acceptance and financing, each tested separately.
\textbf{Long term:} enabling research instruments and productive
infrastructure that leave successors more useful possibilities.

The smallest scientific next step is external mathematical review of
scope, proof and priority. Separately, an authorized sponsor interview
and actual resource quote can test whether the proposed service is worth
buying. No customer contact, minting, trading or public social posting
occurred. This is a research-review edition, distinct from a merged
canonical release.

\textbf{Find your Alpha. Build your SUCCESSOR \(\Omega\).}

\Needspace{7\baselineskip}

\hypertarget{primary-references}{%
\subsection{Primary references}\label{primary-references}}

\hypertarget{refs}{}
\begin{CSLReferences}{1}{0}
\leavevmode\vadjust pre{\hypertarget{ref-emanuel2011}{}}%
Emanuel, Paul Robert. 2011. {``Covering Systems.''} PhD thesis,
Macquarie University. \url{https://doi.org/10.25949/19441907}.

\leavevmode\vadjust pre{\hypertarget{ref-myerson2026}{}}%
Myerson, Gerry. 2026. {``Questions on Finite Congruence Covers, in Open
Problems Which Might Benefit from Computational Experiments.''}
\url{https://mathoverflow.net/a/499017}.

\end{CSLReferences}
\end{document}
